% SPDX-FileCopyrightText: 2026 Will Cook
% SPDX-License-Identifier: CC-BY-4.0
%
% One of the Erdős Problem Notes: a short standalone treatment of a single
% problem in the ErdosProblems expansion library.  Build with
% `tectonic erdos-269-three-prime-running-lcm.tex` or `make`.
\documentclass[11pt]{article}
% SPDX-FileCopyrightText: 2026 Will Cook
% SPDX-License-Identifier: CC-BY-4.0
%
% Shared preamble for the Erdős Problem Notes: the short, standalone,
% problem-owned manuscripts covering the expansion library `ErdosProblems`.
%
% One note owns one Erdős problem.  Each note repeats whatever context its own
% reader needs, so that a reader who arrives at a single problem never has to
% assemble it from the gateway paper.  Deliberate repetition across notes is the
% design, not an oversight.
%
% Source links resolve against one pinned formal-source commit, declared once
% below.  `scripts/check_problem_note_sources.py` verifies that every authored
% (file, line, declaration) triple names that exact declaration in the pinned
% snapshot, so a link cannot silently rot as later work moves lines.

\usepackage{paper-house-style}
\usepackage{array}
\usepackage{booktabs}
\usepackage{enumitem}
\setlist{topsep=0.35em,itemsep=0.2em}
\usepackage[colorlinks=true,linkcolor=linkinternal,urlcolor=linkexternal,
  citecolor=linkinternal]{hyperref}
\hypersetup{bookmarksnumbered=true,bookmarksopen=true,bookmarksopenlevel=1,
  pdfauthor={Will Cook},pdflang={en-GB}}

% ---------------------------------------------------------------------------
% Pinned formal source.  \commit names the checkpoint every link resolves
% against; the checker reads that snapshot rather than the working tree, so the
% printed coordinates stay correct however later waves move declarations.
% ---------------------------------------------------------------------------
\newcommand{\commit}{e5fd26a6d1b1a346430205eab5385b4c9565d004}
\newcommand{\commitshort}{e5fd26a6d1b1}
\newcommand{\repobase}{https://github.com/wcook04/plectis-lean-erdos249-257/blob/\commit}
\newcommand{\sourceurl}{https://github.com/wcook04/plectis-lean-erdos249-257/tree/\commit}
\newcommand{\PX}{ErdosProblems}

% Declaration names stay inside link targets.  The printed label is either a
% spaced, lower-case reading of the declaration name (\lref) or a short authored
% phrase (\lword); neither prints a module path or a raw Lean identifier.
\ExplSyntaxOn
\cs_new_protected:Npn \noteleanlabel #1
  {
    \tl_set:Nx \l_tmpa_tl { \tl_to_str:n {#1} }
    \regex_replace_all:nnN { _+ } { \c{space} } \l_tmpa_tl
    \regex_replace_all:nnN
      { ([a-z0-9])([A-Z]) }
      { \1\c{space}\2 }
      \l_tmpa_tl
    \tl_set:Nx \l_tmpa_tl { \tl_lower_case:n { \tl_use:N \l_tmpa_tl } }
    \tl_use:N \l_tmpa_tl
  }
\ExplSyntaxOff
\newcommand{\leanlabel}[1]{\noteleanlabel{#1}}
\newcommand{\lref}[3]{\href{\repobase/\PX/#1\#L#2}{\leanlabel{#3}}}
\newcommand{\lword}[4]{\href{\repobase/\PX/#1\#L#2}{#4}}
\newcommand{\lloc}[2]{\href{\repobase/\PX/#1\#L#2}{Lean source}}

% Links into a sibling library, named repository-relative.  The expansion
% library is implicit in \lref/\lword, because that is where problem-owned
% work lands.  Several problems have their headline theorem in the older
% reviewed #249/#257 corpus, and a note that could not name that library
% would have to paraphrase a checked statement instead of linking it.
\newcommand{\mref}[3]{\href{\repobase/#1\#L#2}{\leanlabel{#3}}}
\newcommand{\mword}[4]{\href{\repobase/#1\#L#2}{#4}}
\newcommand{\mloc}[2]{\href{\repobase/#1\#L#2}{Lean source}}
\emergencystretch=2em

\newtheorem{theorem}{Theorem}[section]
\newtheorem{proposition}[theorem]{Proposition}
\newtheorem{lemma}[theorem]{Lemma}
\newtheorem{corollary}[theorem]{Corollary}
\theoremstyle{definition}
\newtheorem{definition}[theorem]{Definition}
\newtheorem{example}[theorem]{Example}
\newtheorem{problem}[theorem]{Problem}
\theoremstyle{remark}
\newtheorem*{remark}{Remark}

\newcommand{\N}{\mathbb{N}}
\newcommand{\Npos}{\mathbb{N}_{>0}}
\newcommand{\Z}{\mathbb{Z}}
\newcommand{\Q}{\mathbb{Q}}
\newcommand{\R}{\mathbb{R}}
\newcommand{\lcm}{\operatorname{lcm}}
\newcommand{\ph}{\varphi}

% The series line printed under every note's title block.
\newcommand{\seriesline}{Erd\H{o}s Problem Notes}

% Production provenance.  The wording is fixed by the corpus provenance
% contract and reproduced verbatim in every manuscript in this repository, so
% the papers cannot drift into different accounts of how they were made.
\newcommand{\noteprovenance}{\provenancenote{\emph{Authorship and AI use.} The
prose, formal proofs, and software in this version were drafted and revised by
large-language-model agents under Will Cook's direction.  Cook set the
objectives and acceptance criteria, reviewed and authorised the manuscript, and
is responsible for its contents.  The agents are production tools, not authors.
See \emph{Statements and declarations}.}}

% The status paragraph every note carries, in its own words, before any result.
% Kernel checking establishes the exact Lean proposition; it does not establish
% that the proposition is the interesting one, that it is new, or that the
% Erdős problem is closed.
\newcommand{\statusboundary}{%
  \emph{Status.}  The problem treated here is open, and this note does not
  close it.  Every statement below marked as checked is a proposition that the
  pinned Lean kernel accepts from the sources this note links to, with no
  \texttt{sorry}, no added axiom, and no unchecked evaluation.  That is a claim
  about the formal statement, not about its mathematical interest, its
  novelty, or the original problem.  The unresolved obligations are named
  exactly, in their own section, and none of the finite computations,
  reductions, or no-go results here removes one of them.}

\usepackage{longtable}
\usepackage{colortbl}
% Problem-specific source pin: #269's local-window consumer advances
% independently of the corpus default.
\renewcommand{\commit}{99f4bf47422abbd8757cbb22b50ba079d764d3a7}
\renewcommand{\commitshort}{99f4bf47422a}
\hypersetup{
  pdftitle={No finite separable representation at three prime generators (Erdős Problem 269)},
  pdfsubject={Erdős Problem Notes: Problem 269},
  pdfkeywords={transcendence, Hecke-Mahler values, least common multiple, smooth numbers, lattice sums, Lean 4},
}

% The three pure-power coordinates and the running height.
\newcommand{\hgt}{\operatorname{H}}
\newcommand{\lo}{\lambda}
\newcommand{\hi}{\eta}
\newcommand{\Lprefix}{\operatorname{L}}
\newcommand{\Kernel}{\operatorname{K}}
\newcommand{\authnone}{\textcolor{BrickRed!75!black}{\sffamily\bfseries No result here}}
\newcommand{\authliterature}{\textcolor{BurntOrange!85!black}{\sffamily\bfseries Literature}}
\newcommand{\authlean}{\textcolor{MidnightBlue!82!black}{\sffamily\bfseries Paper + Lean}}
\newcommand{\authexternal}{\textcolor{Violet!80!black}{\sffamily\bfseries Paper + external theorem}}
\newcommand{\authcompute}{\textcolor{Violet!80!black}{\sffamily\bfseries Direct computation}}

\runninghead{Erd\H{o}s \#269: no finite separable representation}
\title{No Finite Separable Representation\\at Three Prime Generators}
\subtitle{Nonsingular minors of every order, two-prime Hecke--Mahler
polynomials, and the open boundary for Erd\H{o}s Problem~\#269}
\author{Will Cook}
\date{July 2026; revised 16 September 2026}

\begin{document}
\maketitle
\noteprovenance

\begin{abstract}
For three distinct primes, the reciprocal running-LCM kernel has nonsingular
minors of every order, with the same indices in every third-coordinate
layer.  Diagonal rescaling leaves a binary floor carry whose threshold
columns determine every finite sampled rank.  This rules out finite exact
separation of the kernel, not rationality of its sum.  For the repeated
$\{2,3,5\}$ series, we give the literal integral tail recurrence, a
quadratic bound, exact denominator clearing and a residue criterion
equivalent to irrationality.  We do not prove the required cofinal escape.
The two-prime comparison reproduces the earlier Hecke--Mahler reduction
attributed to Steve Fan and distinguishes the repeated and distinct-height
values.  A concrete moving-boundary example explains why quadratic
coefficient growth does not by itself yield finite-difference cancellation.
\end{abstract}

\section{Introduction}\label{sec:problem}
Let $P$ be a finite set of at least two primes and let $\mathcal S_P$ be the
positive integers supported on $P$, including $1$.  All sums below retain
multiplicity unless explicitly marked distinct-height.  Put
\[
 \Lprefix(x)=\lcm\{u\in\mathcal S_P:u\le x\},\qquad
 \mathcal R_P=\sum_{u\in\mathcal S_P}\Lprefix(u)^{-1}.
\]
The historical question is whether $\mathcal R_P$ is irrational
\cite[p.~65]{erdosgraham1980}\cite[p.~106]{erdos1988}.
We address the repeated series, not the distinct-height assertion recorded
by Erd\H{o}s in his earlier letter~\cite{erdos1974letter}.
For $P=\{p,q,r\}$, define
\[
 \hgt(x)=p^{\lfloor\log_p x\rfloor}
 q^{\lfloor\log_q x\rfloor}r^{\lfloor\log_r x\rfloor},\qquad
 \Kernel(i,j,k)=\hgt(p^iq^jr^k)^{-1}.
\]
Prime-power divisibility gives $\Lprefix=\hgt$
(Proposition~\ref{res:lcm}).  The supplied forum record credits Steve Fan's
post of 26 June 2026 with this identity for every finite prime set and the
two-prime factorisation, reduction and transcendence conclusion
\cite{fan2026comment}.  We retain that attribution; the live thread was not
accessible during this revision.

With two primes, the corresponding kernel factors into a function of each
exponent separately, reducing the double sum to a product.  The first
question at three primes is whether finitely many such products still
suffice.  Theorem~\ref{res:infinite-rank} rules this out: one floor carry
survives after all row and column factors have been removed.  The proof uses
only this rescaling, the density of the rotations by $\log_r p$ and $\log_r q$
taken separately, and a staircase determinant, and one choice of indices
serves every third-coordinate layer.  We first
explain that obstruction, then recover the two-prime comparison, and finally
identify the arithmetic of the actual three-prime tails.  The last step is
necessary because the rank of a kernel does not determine the irrationality
of its sum.

\paragraph{Literature and scope.}
Fixed-prime semigroups and their neighbouring elements are classical objects;
see Tijdeman--Meijer \cite{tijdemanmeijer1974} and the recent two-prime study
of Languasco, Luca, Moree and Togb\'e \cite{languasco2025}.  No asymptotic
counting theorem from those works is needed for our elementary shell bounds.
Kova\v{c}--Tao \cite{kovactao2024} concerns neighbouring irrationality
problems for distinct unit fractions, not this repeated-height sum.
The rank theorem below concerns equality of kernels with arbitrary
rational-valued separated factors.  It does not exclude a different
analytic representation of the scalar value.

\paragraph{Reading guide.}
Section~\ref{sec:rank} proves the representation theorem.  The two-prime
comparison in Section~\ref{sec:two-prime} uses an external value theorem.
Sections~\ref{sec:actual-orbit} and~\ref{sec:escape} define the actual
arithmetic quantities and one open target, in equivalent tail and window
forms.  Source coordinates are collected in Appendix~\ref{app:sources} and
the accompanying concordance, not used as substitutes for proofs.

\section{The binary carry and arbitrary-order rank}\label{sec:rank}

\begin{theorem}[arbitrary-order non-separability]\label{res:infinite-rank}
Let $p,q,r$ be primes with $p\ne q$, $p\ne r$ and $q\ne r$.  For every
$n\ge0$ there are injective maps $I,J:\{0,\ldots,n-1\}\to\N$ such that, for
every $k\ge0$,
\[
 \det\bigl(\Kernel(I(a),J(b),k)\bigr)_{0\le a,b<n}\ne0.
\]
Consequently, for no finite $d$ do there exist rational-valued functions
$f_\ell(i)$ and $G_\ell(j,k)$, $0\le\ell<d$, satisfying
\[
 \Kernel(i,j,k)=\sum_{\ell<d}f_\ell(i)G_\ell(j,k)
 \qquad\hbox{for all }i,j,k.
\]
\end{theorem}

\begin{proof}
Put $c=r^{-1}$, $x_i=\{i\log_r p\}$ and $y_j=\{j\log_r q\}$.
Splitting the floor exponents gives
\[
 \Kernel(i,j,k)=U_i(k)^{-1}C_{ij}V_j(k)^{-1},\qquad
 C_{ij}=c^{\mathbf1_{\{x_i+y_j\ge1\}}},
\]
where
\[
\begin{aligned}
 U_i(k)&=p^i q^{\lfloor\log_q(p^ir^k)\rfloor}
            r^{k+\lfloor\log_r p^i\rfloor},\\
 V_j(k)&=q^j p^{\lfloor\log_p(q^jr^k)\rfloor}
            r^{\lfloor\log_r q^j\rfloor}.
\end{aligned}
\]
Both factors are positive.  Distinct primes make $\log_r p$ and
$\log_r q$ irrational, so each individual rotation is dense.
No joint density of $(x_i,y_j)$ is needed.  For $n\ge1$, use density of
$(x_i)$ to choose distinct row indices with
$0<x_{I(0)}<\cdots<x_{I(n-1)}<1$.  Now put
$s_b=1-y_{J(b)}$.  Density of $(y_j)$ lets us choose

\[
 s_0\in(0,x_{I(0)}),\qquad
 s_b\in(x_{I(b-1)},x_{I(b)})\quad(1\le b<n).
\]

The intervals for the $s_b$ are disjoint, so the column indices are
distinct.  The strict inequalities also avoid the threshold itself, and give
$x_{I(a)}+y_{J(b)}\ge1$ exactly when $b\le a$.  Thus multiplying row $a$
by $U_{I(a)}(k)$ and column $b$ by $V_{J(b)}(k)$ reduces the selected kernel
minor to
\[
 T_n(c)=\begin{pmatrix}
 c&1&\cdots&1\\
 c&c&\cdots&1\\
 \vdots&\vdots&\ddots&\vdots\\
 c&c&\cdots&c
 \end{pmatrix},\qquad
 \det T_n(c)=c(c-1)^{n-1}.
\]
Indeed, subtracting each preceding row from the next, working upwards from
the last row, leaves diagonal entries $c,c-1,\ldots,c-1$.
Before normalisation, the determinant is this nonzero number times
\[
 \prod_{a<n}U_{I(a)}(k)^{-1}\prod_{b<n}V_{J(b)}(k)^{-1}.
\]
The factors are nonzero, and the same $I,J$ work for every $k$.  The empty
minor for $n=0$ equals $1$.

For the final assertion, suppose a separation with $d$ summands existed and
fix any $k$.  On the rows $I(0),\ldots,I(d)$ and columns
$J(0),\ldots,J(d)$ its matrix would be the product of the $(d+1)\times d$
matrix $(f_\ell(I(a)))_{a,\ell}$ and the $d\times(d+1)$ matrix
$(G_\ell(J(b),k))_{\ell,b}$.  Its rank would be at most $d$, so its
determinant would vanish, contrary to the minor just constructed.
\end{proof}

\paragraph{Scope of the rank obstruction.}
The proof works over $\mathbb R$ or $\mathbb C$ as well as $\mathbb Q$:
a factorisation through a $d$-dimensional space still has rank at most $d$.
The displayed theorem retains the rational-valued scope of the supplied
formal statement.  Also $p\ne q$ is not used by the carry argument itself;
it is retained because the running-LCM application has three distinct
prime generators.

\begin{corollary}[the same minors modulo every admissible denominator]
\label{res:admissible-modular-minors}
For $(p,q,r)=(2,3,5)$ and every $n\ge1$, there are injective maps
$I,J:\{0,\ldots,n-1\}\to\N$, chosen independently of $B$ and $k$, such that
for every $B\ge2$ coprime to $30$ and every $k\ge0$, the selected $n\times n$
kernel matrix has unit determinant over $\mathbb Z/B\mathbb Z$, with each
reciprocal prime power interpreted by its modular inverse.
\end{corollary}
\begin{proof}
Choose the maps from Theorem~\ref{res:infinite-rank}.  Every row and column
factor is a unit modulo $B$.  The normalised determinant is
$5^{-1}(-4/5)^{n-1}$, also a unit.
\end{proof}





\begin{example}\label{res:rank}
For $(p,q,r)=(2,3,5)$, the leading two-by-two determinant is $-1/15$:
its four entries are $1,1/6,1/2,1/60$, so it equals $1/60-1/12$.
\end{example}



\paragraph{The indices must be selected.}
The leading $4\times4$ block at $\{2,3,5\}$ is singular:
\[
 \Kernel(3,j,0)=\frac1{120}\Kernel(0,j,0)\qquad(0\le j<4).
\]
These equalities follow by evaluating the height at $3^j$ and $8\cdot3^j$.
At $j=4$ the difference is $-1/19440000$.
Thus the leading minors do not supply the arbitrary-order theorem.


\begin{proposition}[finite sampled cut rank]\label{res:finite-cut-rank}
Let $m\ge1$ and let $c$ lie in a field with $c\ne 0,1$.
For $0\le k\le m$ write $v_k$ for the length-$m$ column with a $1$ in each of
the first $k$ coordinates and $c$ thereafter.
A matrix whose distinct columns are $v_k$ for $k$ in a nonempty set $E$ has
rank $|E|-\mathbf 1_{\{0,m\}\subseteq E}$.
\end{proposition}

\begin{proof}
List $E=\{k_1<\cdots<k_t\}$.  The $t-1$ differences
$v_{k_{j+1}}-v_{k_j}=(1-c)\mathbf 1_{\{k_j,\ldots,k_{j+1}-1\}}$
have disjoint nonempty supports, hence are linearly independent.
If $k_1>0$ or $k_t<m$, their union misses a coordinate where $v_{k_1}$ is
nonzero, so the rank is $t$.
If $k_1=0$ and $k_t=m$, then $v_0=c\,\mathbf 1$ lies in the span of the
differences (their sum is $(1-c)\mathbf 1$), so the rank is $t-1$.
\end{proof}



After ordering the sampled row phases, the columns have exactly the forms
$v_k$ above; repeated rows or columns add no rank.  Thus the formula
classifies every nonempty finite sample after diagonal rescaling, uniformly
in the third coordinate.  It also gives an exact algorithm: order the
rationals $p^i/r^{\lfloor\log_r p^i\rfloor}$, count those strictly below
$r^{\lfloor\log_r q^j\rfloor+1}/q^j$, and apply the endpoint correction.
Integer comparisons suffice; no floating-point decision at a threshold is
needed.


\paragraph{The norm matters.}
For the infinite normalised carry matrix,
\[
 \inf_{F\text{ of finite separated rank}}\|C-F\|_\infty=(1-c)/2.
\]
For $c=1/r$, distinct columns are exactly $1-c$ apart in the supremum
norm.  An error $\varepsilon<(1-c)/2$ would place the approximating columns
in a bounded subset of their finite-dimensional span, with pairwise distance
at least $1-c-2\varepsilon>0$.  Such a subset is totally bounded, a
contradiction.  This argument uses bounded columns, not bounded individual
factors in a separated expression.  The constant $(1+c)/2$ attains the bound.

Finite restrictions have a different behaviour. At $c=1/5$,
\[
 T=\begin{pmatrix}1/5&1\\1/5&1/5\end{pmatrix},\qquad
 A=\begin{pmatrix}3/10&9/10\\1/10&3/10\end{pmatrix}
\]
satisfy $\det A=0$ and $\|T-A\|_{\max}=1/10<2/5$.
For the original three-index kernel, put
$K^{(N)}(i,j,k)=\mathbf1_{\{i<N\}}K(i,j,k)$.  This has separated rank at
most $N$ across $i\mid(j,k)$, and the $\ell^1(\mathbb N^3)$ error satisfies
\[
 \|K-K^{(N)}\|_{\ell^1}
 \le\frac{pqr\,p^{-3N}}{(1-p^{-3})(1-q^{-3})(1-r^{-3})}.
\]
This follows from $\hgt(x)>x^3/(pqr)$. Thus the infinite uniform obstruction
coexists with geometric approximation in the summation norm. The scalar
irrationality question requires arithmetic information about the actual
multiplicities.

\section{The two-prime comparison}\label{sec:two-prime}
For distinct primes $p<q$, write
$L_{p,q}(t)=p^{\lfloor\log_p t\rfloor}q^{\lfloor\log_q t\rfloor}$.
Prime-power divisibility identifies this product with the running LCM.
Let $\mathcal R_{p,q}$ retain every smooth-number multiplicity, and let
$\mathcal D_{p,q}$ retain the initial value and one reciprocal at each
positive pure-power jump.  Fan's forum post gives the repeated two-prime
factorisation and its quadratic reduction to a single Hecke--Mahler value, and
the same boundary-channel calculation yields the de-duplicated affine formula
displayed below~\cite{fan2026comment}.  We give the calculation to fix
normalisations and to distinguish repeated from distinct-height sums.

\begin{theorem}[both two-prime sums]
\label{res:two-prime-transcendence}
\label{res:two-prime-repeated-transcendence}
Let $p<q$ be distinct primes.  Put $\theta=\log p/\log q$ and
$A=\sum_{n\ge0}p^{-n}q^{-\lfloor n\theta\rfloor}$.  Let
$\mathcal R_{p,q}$ sum the reciprocal running LCM at every positive
$\{p,q\}$-smooth integer, and let $\mathcal D_{p,q}$ count each distinct
running LCM once.  Then
\begin{equation}\label{eq:two-prime-affine}
 \mathcal D_{p,q}=\frac{(q-p)A+p}{q-1},\qquad
 \mathcal R_{p,q}=\frac{(p+q-1)A-(p-1)A^2}{q-1}.
\end{equation}
Both numbers are transcendental.
\end{theorem}
\begin{proof}
Set $x=1/p$, $y=1/q$, $m_n=\lfloor n\theta\rfloor$ and
$\delta_n=m_{n+1}-m_n$.  Unique factorisation makes $\theta$ irrational,
and $0<\theta<1$ gives $\delta_n\in\{0,1\}$.  The initial value and the
$p$-channel contribute $A=\sum_{n\ge0}x^ny^{m_n}$.  There is one $q$-power
between $p^n$ and $p^{n+1}$ precisely when $\delta_n=1$, so the other
channel contributes
\[
 B=\sum_{n\ge0}\delta_nx^ny^{m_n+1},\qquad \mathcal D_{p,q}=A+B.
\]
All these series converge absolutely.  Since
$y^{m_{n+1}}-y^{m_n}=\delta_ny^{m_n}(y-1)$, shifting the series for $A$
gives $A-1-xA=x(y-1)B/y$.  Therefore
\[
 B=\frac{p-(p-1)A}{q-1}.
\]
At $p^iq^j$ the height is
$p^{i+\lfloor j/\theta\rfloor}q^{j+m_i}$, hence
\[
 \mathcal R_{p,q}
 =A\sum_{j\ge0}y^jx^{\lfloor j/\theta\rfloor}=A(1+B).
\]
For the last equality, each $j\ge1$ corresponds to
$n=\lfloor j/\theta\rfloor$ with $m_n=j-1$ and $\delta_n=1$.
These identities prove~\eqref{eq:two-prime-affine}.

For the Hecke--Mahler series
$F_\theta(x,y)=\sum_{n\ge1}\sum_{k=1}^{m_n}x^ny^k$, geometric summation gives
\begin{equation}\label{eq:hecke-mahler-boundary}
 A=\frac1{1-x}-\frac{1-y}{y}F_\theta(x,y).
\end{equation}
Bugeaud and Laurent's theorem \cite[Theorem~1.1]{bugeaudlaurent2023} applies
with intercept $\rho=0$, $\beta=x$ and $\alpha=y$: the slope $\theta$ is
irrational and lies in $(0,1)$, $x$ and $y$ are nonzero algebraic numbers,
$|x|<1$ and $|xy^\theta|=p^{-2}<1$.  Hence $F_\theta(x,y)$ is transcendental;
this case $\rho=0$ is due to Loxton and van der Poorten
\cite[Theorem~8, p.~40]{loxtonvdp1977}.  Thus $A$ is transcendental.
The displayed affine and quadratic polynomials are nonconstant, so an
algebraic value of either would force $A$ to be algebraic.
\end{proof}

The two-prime reduction expresses both sums through one value.  The third
prime leaves the binary carry of Section~\ref{sec:rank}.  Its rank theorem
concerns that function; the scalar arithmetic depends on the multiplicities
introduced next.  The de-duplicated irrationality assertion is historical
\cite{erdos1974letter}.

\section{Exact multiplicities and normalised tails}
\label{sec:lcm}\label{sec:cells}\label{sec:fibre}\label{sec:shell}
\label{sec:actual-orbit}

\begin{proposition}[the running least common multiple]\label{res:lcm}
Let $p,q,r$ be pairwise distinct primes and $x\ge1$.  Then the running LCM equals the three-prime height:
$\Lprefix(x)=\hgt(x)$.
\end{proposition}
\begin{proof}
Every smooth $n\le x$ has prime exponents bounded by the corresponding
integer logarithms, so $n\mid\hgt(x)$.  Conversely the three maximal pure
powers occur among those smooth numbers; their product divides the running
LCM because they are pairwise coprime.
\end{proof}

The inequalities $x/p<p^{\lfloor\log_p x\rfloor}\le x$ for each of the
three primes also give
\begin{equation}\label{res:cube}
 x^3/(pqr)<\hgt(x)\le x^3.
\end{equation}

\begin{proposition}[cells and jumps]\label{res:cell}
The running LCM is constant when the three integer logarithms are constant.
A jump in exactly one logarithm multiplies it by the corresponding prime.
The first $n$ positive powers in each channel, together with $1$, form
$3n+1$ distinct points.
\end{proposition}
\begin{proof}
The first two claims follow from the height formula.  Positive powers from
different prime channels cannot coincide, by unique factorisation, and
none is $1$.
\end{proof}

\begin{proposition}[height-fibre regrouping]\label{res:fibre-prop}
For a finite exponent box $\mathcal B$, set
$F(H)=\{(i,j,k)\in\mathcal B:\hgt(p^iq^jr^k)=H\}$.  Then
\begin{equation}\label{res:fibre}
 \sum_{(i,j,k)\in\mathcal B}\Kernel(i,j,k)=\sum_H\frac{\#F(H)}H.
\end{equation}
\end{proposition}
\begin{proof}
Each term in $F(H)$ equals $1/H$.  The multiplicities remain present in
every subsequent infinite sum.
\end{proof}



From now on let $P=\{2,3,5\}$ and $S=\mathcal R_P$.  For $a\ge0$ define
\[
 s_a=\sum_{\substack{x\text{ smooth}\\2^a\le x<2^{a+1}}}\frac1{\hgt(x)},
 \quad T_a=\sum_{j\ge0}s_{a+j},\quad
 h_a=\frac{\hgt(2^a)}2,\quad X_a=h_aT_a.
\]
Thus $T_a$ is the raw tail, $X_a$ its normalised state, and
$P_a:=\hgt(2^a)=2h_a$.  In particular $h_0=1/2$, while $h_a$ is a
positive integer for $a\ge1$.  The literal radix and forcing numerator are
\begin{equation}\label{eq:actual-digit}
 b_a=\frac{\hgt(2^{a+1})}{\hgt(2^a)},\qquad
 m_a=\sum_{\substack{x\text{ smooth}\\2^a\le x<2^{a+1}}}
       \frac{\hgt(2^{a+1})}{2\hgt(x)}.
\end{equation}

\begin{lemma}[integer forcing and the four radices]
\label{res:dyadic-alphabet}
For every $a\ge0$, $m_a$ is a positive integer and
$b_a\in\{2,6,10,30\}$.  The word ``numerator'' does not impose the
positional-digit restriction $m_a<b_a$; that restriction need not hold.
\end{lemma}
\begin{proof}
If $x<2^{a+1}$, the two-exponent of $\hgt(x)$ is at most $a$, while the
other exponents are bounded by those of $\hgt(2^{a+1})$.  Hence
$2\hgt(x)\mid\hgt(2^{a+1})$.  Each forcing summand is an integer,
and the shell contains $2^a$.  Between consecutive two-powers there is at most one three-power and at
most one five-power: two successive powers in either channel have ratio
strictly greater than two.  Their optional
factors, together with the terminal factor two, give the four radices.
\end{proof}



Since $h_{a+1}=b_0b_1\cdots b_a/2$, the value
$S/2=\sum_{a\ge0}m_a/(b_0b_1\cdots b_a)$ is a Cantor series and
$X_a=\sum_{j\ge a}m_j/(b_a\cdots b_j)$ is its normalised tail.  Such tails,
with their integer carry recurrence, are the classical tool for the
rationality of Cantor series.  Erd\H{o}s and Straus characterise rationality
by integer carries \cite[Theorem~2.1 and (2.4), pp.~85--86]{erdosstraus1974},
and Han\v{c}l and Tijdeman give a practical form of that criterion through the
tails \cite[\S\S2--3 and Theorem~3.1, pp.~372--375]{hancltijdeman2004}.  Both
criteria assume that the $n$th numerator is $o(a_{n-1}a_n)$, where $a_n$ is
the $n$th base.  Here $m_a\ge1$ and $b_{a-1}b_a\le900$, so the recurrence,
the bound and the denominator clearing below are proved directly for this
series.
The polynomial criterion of Han\v{c}l--Tijdeman
\cite[Theorems~2.2 and 3.1]{hancltijdeman2008} assumes a nonconstant
polynomial radix and does not cover this bounded sequence.  Their
Theorem~4.2 is a separate, more general telescoping principle: it requires
an exact finite-product decomposition and a small transformed numerator.
Neither a quadratic bound nor a polynomial carry ansatz supplies those
hypotheses.  The prior Isabelle/HOL formalisation of the classical
Erd\H{o}s--Straus criteria is due to Koutsoukou-Argyraki and Li
\cite{afperdosstraus2020}; it is not a verification of the present series.

\begin{proposition}[the actual recurrence and a polynomial bound]
\label{res:actual-orbit}
The series defining $S,T_a$ converge.  For every $a\ge0$,
\[
 X_{a+1}=b_aX_a-m_a,\qquad
 0<X_a\le\frac{8640}{343}(a+1)^2<90(a+1)^2.
\]
For every integer $B\ge1$, either some $BX_a$ is integral and all later
states are integral, or $\operatorname{dist}(BX_a,\mathbb Z)\ge1/31$
at arbitrarily large indices.
\end{proposition}
\begin{proof}
A dyadic shell contains at most one two-exponent for each pair of three-
and five-exponents.  Each of the latter lies between $0$ and $a$, so there are at most
$(a+1)^2$ terms.  By~\eqref{res:cube},
$s_a\le30(a+1)^2/8^a$.  Therefore
\[
 X_a\le15(a+1)^2\sum_{j\ge0}\frac{(j+1)^2}{8^j}
      =\frac{8640}{343}(a+1)^2.
\]
This proves convergence and the bound.  Splitting the first shell gives
the recurrence because $m_a=h_{a+1}s_a$ and $h_{a+1}=b_ah_a$.
Put $Y_a=BX_a$.  If all sufficiently late distances are strictly below
$1/31$, write $Y_a=z_a+e_a$ with $z_a\in\mathbb Z$ and $|e_a|<1/31$.
Then $e_{a+1}-b_ae_a$ is an integer of absolute value less than one, so
$e_{a+1}=b_ae_a$.  The factors $b_a\ge2$ force a bounded such error to be
zero.  An integral state propagates by the integer recurrence.
\end{proof}

The window criterion needs only a valid carry cap; $90B(a+1)^2$ suffices.
The long record proves $X_a\le Q(n_a)$ and the smaller bound
$X_a\le\widetilde Q(n_a)$, where
$n_a=a+\lfloor\log_3 2^a\rfloor+\lfloor\log_5 2^a\rfloor$.
It keeps these sharper caps separate from the coarse cap used here.

\begin{theorem}[actual rationality-to-reduced-carry bridge]
\label{res:denominator-reduction}
If $S=A/D$ in lowest terms, where
$D=2^u3^v5^wB$ and $\gcd(B,30)=1$, then for every
$a\ge a_0=u+1+2v+3w$,
\[
 d_a=BX_a\in\mathbb Z_{>0},\qquad
 d_{a+1}=b_ad_a-Bm_a,\qquad d_a\le90B(a+1)^2.
\]
\end{theorem}
\begin{proof}
For $a\ge1$, strict boundary clearing gives
\begin{equation}\label{eq:prefix-lattice}
 X_a=h_aS-Z_a,\qquad
 Z_a=\sum_{\substack{x\text{ smooth}\\x<2^a}}\frac{h_a}{\hgt(x)}\in\mathbb Z.
\end{equation}
The two-exponent of $h_a$ is $a-1$.  Moreover $a\ge2v$ gives
$2^a\ge3^v$, and $a\ge3w$ gives $2^a\ge5^w$.
Thus $2^u3^v5^w\mid h_a$ after the stated onset, and
$BX_a=h_aA/(2^u3^v5^w)-BZ_a$ is integral.  Positivity, the recurrence
and the bound follow from Proposition~\ref{res:actual-orbit}.
\end{proof}



\begin{corollary}[the exact clearing onset]\label{res:exact-onset}
Under the same lowest-terms hypothesis, put $M=2^u3^v5^w$ and let
$\operatorname{den}$ denote the positive reduced denominator.  For $a\ge1$,
\[
 \operatorname{den}(BX_a)=\frac{M}{\gcd(M,h_a)}.
\]
Consequently $BX_a$ is integral exactly when
$2^a\ge\max(2^{u+1},3^v,5^w)$.  The first such $a$ can be found by
integer comparisons, without logarithmic rounding.
\end{corollary}
\begin{proof}
Equation~\eqref{eq:prefix-lattice} gives
$BX_a=h_aA/M-BZ_a$.  Since $\gcd(A,M)=1$, reducing that fraction gives
the displayed denominator.  The three divisibility conditions for $h_a$
are precisely the three inequalities stated above.
\end{proof}


\section{A window test and the remaining arithmetic}
\label{sec:escape}\label{sec:open}
The actual digits in~\eqref{eq:actual-digit} define
\[
 W_{\ell,0}=1,\quad F_{\ell,0}=0,\qquad
 W_{\ell,h+1}=b_{\ell+h}W_{\ell,h},\quad
 F_{\ell,h+1}=b_{\ell+h}F_{\ell,h}+m_{\ell+h}.
\]
Induction on $h$ gives
\begin{equation}\label{eq:actual-tail}
 X_{\ell+h}=W_{\ell,h}X_\ell-F_{\ell,h}.
\end{equation}
For integers $W\ge1$ and $t$, define
$\operatorname{lpr}_W(t)=1+((t-1)\bmod W)$, using a remainder in
$\{0,\ldots,W-1\}$.  Thus a zero residue is represented by $W$, not $0$.
Throughout this section use the valid cap
\begin{equation}\label{eq:actual-bound}
 K(B,a)=90B(a+1)^2.
\end{equation}

\begin{lemma}[finite endpoint obstruction]\label{res:consumer}
If $d$ is a positive integer with $d\le K$ and $d\equiv -BF\pmod W$,
where $W\ge1$, then $\operatorname{lpr}_W(-BF)\le K$.
\end{lemma}
\begin{proof}
Every positive representative of the residue is at least its least
positive representative, so $\operatorname{lpr}_W(-BF)\le d\le K$.
\end{proof}



\begin{theorem}[the actual window criterion]\label{res:windowconsumer}
The number $S$ is irrational if and only if
\begin{equation}\label{eq:escape}
 \begin{gathered}
 \text{for every }B\ge1\text{ with }\gcd(B,30)=1
 \text{ and every }a_0\ge1,\\
 \text{there are }\ell\ge a_0,\ h\ge1\text{ such that}
 \operatorname{lpr}_{W_{\ell,h}}(-BF_{\ell,h})>K(B,\ell+h).
 \end{gathered}
\end{equation}
\end{theorem}
\begin{proof}
If $S$ is rational, Theorem~\ref{res:denominator-reduction} supplies a
positive integral carry $d_a=BX_a$ after its onset.  Choose a window in
\eqref{eq:escape} beyond that onset.  Equation~\eqref{eq:actual-tail}
gives $d_{\ell+h}\equiv-BF_{\ell,h}\pmod{W_{\ell,h}}$, contradicting
Lemma~\ref{res:consumer} and the valid cap.

Conversely, let $S$ be irrational and fix $B,\ell\ge1$.
Equation~\eqref{eq:prefix-lattice} makes $BX_\ell$ nonintegral.
Put $\delta=\lceil BX_\ell\rceil-BX_\ell\in(0,1)$.
Since $W_{\ell,h}\ge2^h$ and $X_{\ell+h}=O((\ell+h+1)^2)$, for all
sufficiently large $h$ the number
$\delta+BX_{\ell+h}/W_{\ell,h}$ lies in $(0,1)$.  The window identity
then gives the exact equality
\[
 \operatorname{lpr}_{W_{\ell,h}}(-BF_{\ell,h})
 =\delta W_{\ell,h}+BX_{\ell+h}.
\]
Its exponentially growing first term eventually exceeds
$K(B,\ell+h)$.  This proves~\eqref{eq:escape}, in fact from each fixed start.
\end{proof}



The validity of the cap is needed in the rational direction; an upper
polynomial growth condition by itself would not suffice.  The zero cap,
for example, makes positive-residue escape automatic.
The criterion identifies the target in window coordinates.  The unresolved
point is the source-specific exclusion below.

\begin{problem}[the repeated three-prime target]
Prove that $S=\mathcal R_{\{2,3,5\}}$ is irrational.
\end{problem}
\begin{problem}[exact nonintegrality of every reduced tail]\label{prob:tails269}
Prove that for every $a\ge1$ and every integer $B\ge1$ coprime to $30$,
\begin{equation}\label{eq:tail-nonintegrality}
 BX_a\notin\mathbb Z.
\end{equation}
\end{problem}
\begin{problem}[actual cofinal local-window escape]\label{prob:producer}
Prove the cofinal window condition~\eqref{eq:escape}.
\end{problem}
Indeed, one integral scaled tail makes $S$ rational by
\eqref{eq:prefix-lattice}; conversely, a rational $S$ has such tails after
denominator clearing.  These are equivalent formulations of one target,
not two independent missing hypotheses.

\paragraph{A coefficient obstruction, not a transferred theorem.}
Luca, Ouaknine and Worrell
\cite[Definition~5, Theorems~6 and 8, Claim~10]{lucaouaknineworrell2025}
use finite-difference cancellation away from floor crossings, followed by
expanding-gap and relative-variation estimates.  For our literal numerators,
put $\gamma_{a,t}=P_tP_{a+1}/P_{a+t+1}$ and
$D_{r,a}=15\sum_{j=0}^3(-1)^j\binom3j\gamma_{a,jr}m_{a+jr}$.
At $a=0,r=35$, exact comparisons give
$\lfloor\log_3 2^{35\nu}\rfloor=22\nu$,
$\lfloor\log_5 2^{35\nu}\rfloor=15\nu$ and $b_{35\nu}=2$ for
$0\le\nu\le3$.  These facts together eliminate the relevant floor-addition
crossings and make all four corrections $1$.
Nevertheless
\[
 D_{35,0}=15(1-3\cdot195+3\cdot723-1582)=45\ne0.
\]
The long record proves this by exact lattice counts and isolates the
uncancelled boundary strips.  It refutes automatic cubic cancellation,
not every possible operator or an eventual sparse-defect construction.

A fixed-base expansion is not impossible: the direct base-$30$ recoding
has coefficients $e_a=m_a30^{a+1}/P_{a+1}\in\mathbb Z_{>0}$, but
$e_a\ge(15/4)^{a+1}$, so it fails the polynomial-growth hypothesis of the
cited criterion.  The finite-alphabet echoing criterion
\cite{kebis2024echoing} and the base-expansion complexity theorem
\cite[Theorem~1]{adamczewskibugeaud2007} likewise require properties of a
value-preserving coefficient or digit sequence, not merely of $(b_a)$.

\paragraph{A separate distinct-height programme.}
For $\mathcal D_{2,3,5}$, recovering a proof of the historical irrationality
assertion or establishing transcendence is a different task.  An analytic
approach must first construct a specified function and an exact value
identity, then verify a published value theorem's hypotheses.  In
particular, the multivariate Mahler theory of Adamczewski and Faverjon
\cite{adamczewskifaverjon2026} requires a functional system and suitable
evaluation data; having a two-dimensional lattice count is not enough.
No such system is established here.

\paragraph{The information the remaining argument must use.}
The bounded-radix alternative permits the integral branch.  If $(g_a)$ is an
unbounded sequence of positive integers with bounded successive ratios, the
real $\xi$ with $\operatorname{dist}(g_a\xi,\mathbb Z)\to0$ form a countable
set.  For increasing sequences this is due to Erd\H{o}s and Taylor
\cite[Theorem~1, p.~600]{erdostaylor1957}, who obtain it by modifying
Eggleston's proof of his Theorem~16, and Fan proves the form stated here
\cite[Lemma~3.1, p.~7]{fan2026strongly}.  For $g_a=Bh_a$ $(a\ge1)$ the
ratios are the integers $b_a\in\{2,6,10,30\}$, and the argument of
Proposition~\ref{res:actual-orbit} identifies the set: it consists of the
$\xi$ for which $Bh_a\xi$ is eventually an integer.
By~\eqref{eq:prefix-lattice} and Theorem~\ref{res:denominator-reduction} it
contains $S$ for some $B$ exactly when $S$ is rational.  Infinite rank
and the modular minors of Corollary~\ref{res:admissible-modular-minors}
constrain the kernel, while a scalar argument must also use its exact
multiplicities.  The long record gives further conditional rigidity statements, but an
argument for $S$ must verify their block identities for the actual weighted
multiplicities.  Neither a finite-dimensional ansatz nor a countability
statement does that.

\pdfbookmark[1]{Statements and declarations}{statements-declarations}
\subsection*{Statements and declarations}
\paragraph{Proof sources.}
The two-prime deduction uses an external transcendence theorem and is not
formalised here.  The arbitrary-order rank theorem and literal tail bridge
have the named Lean sources recorded in the concordance.  Two verification tracks must be
kept separate: the public repository and the selected Palomar release.
The supplied 16 September 2026 inventory records a successful public CI build
at commit \nolinkurl{99f4bf47422abbd8757cbb22b50ba079d764d3a7}.
The separate Palomar release at \texttt{52f29ad1} selects the
arbitrary-order uniform-minor and non-separation theorems themselves, not
merely the $2\times2$ fixture.  Its recorded Comparator replay therefore
has that broader scope.  These are attached verification receipts, not a
fresh build performed for this editorial revision, and do not certify every
sentence of either paper or the irrationality of $S$.
The elementary modular-minor corollary is proved above.  Its older linked
module is not among the attached source files, so its current build status
is not inferred from the inventory.  The exact-denominator corollary and the boundary example added in this
revision have ordinary proofs and exact checks, but no new formalisation
claim.  The required cofinal escape remains unproved.

\paragraph{Artefact and data availability.}
Each source link is pinned to its own immutable revision.  The exact-reference
inventory records those revisions and their bytewise comparison with the
public source commit recorded in the supplied inventory.
The extended reasoning record and the accompanying source retain the
detailed shell bounds, finite experiments and alternative representations
with their exact hypotheses.  Supplied later-source labels and public links remain separate
source coordinates.

\paragraph{Funding and competing interests.} This work received no external
funding.  The author declares no competing interests.
\paragraph{Acknowledgements.} The problem numbering and the historical
status snapshot are taken from the Erd\H{o}s Problems catalogue maintained
by Thomas Bloom~\cite{erdosproblems}.  Earlier mathematical inputs are
credited at their points of use.

\appendix
\section{Guide to the formal sources}\label{app:sources}
The mathematical dependency chain is: height identity, kernel rescaling,
threshold determinant; and, independently, literal shell forcing, tail
recurrence, denominator clearing and residue contradiction.  The historical declaration inventory is retained in the long record and
the machine-readable source index.  The package also includes an exact
concordance of the links removed from the main narrative.  Moving a link
changes its location, not its immutable revision or evidence status.

\begin{center}\small
\begin{tabular}{p{0.27\linewidth}p{0.64\linewidth}}
\hline
Statement family & Source and scope\\ \hline
Kernel rank & \texttt{KernelCarryRank}: arbitrary-order uniform minors and
no finite rational separation; also selected in the Palomar release.\\
Literal series & \texttt{DyadicShellSummability},
\texttt{RationalityCarryBridge}: actual summable tails and denominator
clearing, not merely an abstract recurrence.\\
Quadratic bounds & \texttt{PaperR8RankMajorant} and
\texttt{ActualSharpTailMajorantR10}: the actual $Q$ bound and the smaller
jump-constrained bound recorded in the long paper.\\
Residue equivalence & \texttt{CofinalWindowEscapeEquivalence} and
\texttt{R12/OcticWindowBand}: conditional consumer and exact equivalence,
not an unconditional cofinal producer.\\
External input & Bugeaud--Laurent and Loxton--van der Poorten for the
two-prime theorem; no Lean value theorem claimed.\\ \hline
\end{tabular}
\end{center}
The public inventory supplied for this revision pins commit
\nolinkurl{99f4bf47422abbd8757cbb22b50ba079d764d3a7}; historical hyperlinks in
the text retain their own revisions.  A source location, a CI build, a
selected Comparator statement and an axiom audit are different evidence
items.  None may stand in for the others without its own receipt.

% BEGIN pin-faithful-source-inventory
\section{Pinned Lean sources}
\label{sec:pinned-lean-sources}
The following declarations are this note's pin-faithful source inventory.
Line numbers are those of the commit named by \texttt{\textbackslash commit}.
\begin{itemize}
\item \lref{Erdos269/ThreePrimeRunningLcm.lean}{123}{smoothPrefixLcm_eq_threePrimeHeight}
\item \lref{Erdos269/ThreePrimeRunningLcm.lean}{179}{smoothPrefixLcm_eq_of_sameLogCell}
\item \lref{Erdos269/ThreePrimeRunningLcm.lean}{326}{smoothPrefixLcm_firstLogStep}
\item \lref{Erdos269/ThreePrimeRunningLcm.lean}{339}{smoothPrefixLcm_secondLogStep}
\item \lref{Erdos269/ThreePrimeRunningLcm.lean}{352}{smoothPrefixLcm_thirdLogStep}
\item \lref{Erdos269/ThreePrimeRunningLcm.lean}{278}{threePrimeJumpSetWithOrigin_card}
\item \lref{Erdos269/ThreePrimeRunningLcm.lean}{406}{finiteSmoothKernelSum_groupedByHeight}
\item \lref{Erdos269/ThreePrimeRunningLcm.lean}{699}{dyadicBlockBase235_cases}
\item \lref{Erdos269/ThreePrimeRunningLcm.lean}{668}{dyadicInternalPower_exponent_unique}
\item \lref{Erdos269/ThreePrimeRunningLcm.lean}{646}{smoothExponentShell_card_quadratic}
\item \lref{Erdos269/ThreePrimeRunningLcm.lean}{31}{smooth3Val}
\item \lref{Erdos269/ThreePrimeRunningLcm.lean}{36}{threePrimeHeight}
\item \lref{Erdos269/ThreePrimeRunningLcm.lean}{40}{threePrimeKernelQ}
\item \lref{Erdos269/ThreePrimeRunningLcm.lean}{46}{smoothPrefixExponents}
\item \lref{Erdos269/ThreePrimeRunningLcm.lean}{53}{smoothPrefixLcm}
\item \lref{Erdos269/ThreePrimeRunningLcm.lean}{59}{smooth3Val_dvd_threePrimeHeight_of_mem}
\item \lref{Erdos269/ThreePrimeRunningLcm.lean}{77}{smoothPrefixLcm_dvd_threePrimeHeight}
\item \lref{Erdos269/ThreePrimeRunningLcm.lean}{84}{pureFirst_mem_smoothPrefixExponents}
\item \lref{Erdos269/ThreePrimeRunningLcm.lean}{97}{pureSecond_mem_smoothPrefixExponents}
\item \lref{Erdos269/ThreePrimeRunningLcm.lean}{110}{pureThird_mem_smoothPrefixExponents}
\item \lref{Erdos269/ThreePrimeRunningLcm.lean}{163}{SameThreePrimeLogCell}
\item \lref{Erdos269/ThreePrimeRunningLcm.lean}{169}{threePrimeHeight_eq_of_sameLogCell}
\item \lref{Erdos269/ThreePrimeRunningLcm.lean}{191}{threePrimeKernelQ_eq_of_sameLogCell}
\item \lref{Erdos269/ThreePrimeRunningLcm.lean}{204}{positivePrimePowers}
\item \lref{Erdos269/ThreePrimeRunningLcm.lean}{208}{positivePrimePowers_card}
\item \lref{Erdos269/ThreePrimeRunningLcm.lean}{219}{one_not_mem_positivePrimePowers}
\item \lref{Erdos269/ThreePrimeRunningLcm.lean}{229}{positivePrimePowers_disjoint}
\item \lref{Erdos269/ThreePrimeRunningLcm.lean}{243}{threePrimePositiveJumpSet}
\item \lref{Erdos269/ThreePrimeRunningLcm.lean}{249}{threePrimePositiveJumpSet_card}
\item \lref{Erdos269/ThreePrimeRunningLcm.lean}{273}{threePrimeJumpSetWithOrigin}
\item \lref{Erdos269/ThreePrimeRunningLcm.lean}{294}{threePrimeHeight_firstLogStep}
\item \lref{Erdos269/ThreePrimeRunningLcm.lean}{305}{threePrimeHeight_secondLogStep}
\item \lref{Erdos269/ThreePrimeRunningLcm.lean}{316}{threePrimeHeight_thirdLogStep}
\item \lref{Erdos269/ThreePrimeRunningLcm.lean}{368}{smoothExponentBox}
\item \lref{Erdos269/ThreePrimeRunningLcm.lean}{373}{smoothPointHeight}
\item \lref{Erdos269/ThreePrimeRunningLcm.lean}{377}{smoothHeightFiber}
\item \lref{Erdos269/ThreePrimeRunningLcm.lean}{384}{smoothHeightFiber_kernel_sum}
\item \lref{Erdos269/ThreePrimeRunningLcm.lean}{430}{threePrimeHeight_le_cube}
\item \lref{Erdos269/ThreePrimeRunningLcm.lean}{443}{kernel_235_origin}
\item \lref{Erdos269/ThreePrimeRunningLcm.lean}{450}{kernel_235_two}
\item \lref{Erdos269/ThreePrimeRunningLcm.lean}{458}{kernel_235_three}
\item \lref{Erdos269/ThreePrimeRunningLcm.lean}{467}{kernel_235_six}
\item \lref{Erdos269/ThreePrimeRunningLcm.lean}{479}{kernel_235_not_rankOne}
\item \lref{Erdos269/ThreePrimeRunningLcm.lean}{487}{tailStateStep}
\item \lref{Erdos269/ThreePrimeRunningLcm.lean}{490}{tailStateStep_eq}
\item \lref{Erdos269/ThreePrimeRunningLcm.lean}{500}{smoothExponentShell}
\item \lref{Erdos269/ThreePrimeRunningLcm.lean}{509}{exponent_unique_in_short_interval}
\item \lref{Erdos269/ThreePrimeRunningLcm.lean}{543}{smoothExponentShell_card_le_dropFirst}
\item \lref{Erdos269/ThreePrimeRunningLcm.lean}{585}{smoothExponentShell_card_le_dropThird}
\item \lref{Erdos269/ThreePrimeRunningLcm.lean}{629}{sorted_pair_quadratic}
\item \lref{Erdos269/ThreePrimeRunningLcm.lean}{662}{DyadicInternalPower}
\item \lref{Erdos269/ThreePrimeRunningLcm.lean}{711}{dyadicBlockBase235_mem_interval}
\item \lref{Erdos269/ResidueEscape.lean}{26}{leastPositiveResidue}
\item \lref{Erdos269/ResidueEscape.lean}{31}{leastPositiveResidue_pos_le}
\item \lref{Erdos269/ResidueEscape.lean}{52}{leastPositiveResidue_modEq}
\item \lref{Erdos269/ResidueEscape.lean}{71}{ResidueEscapesWindow}
\item \lref{Erdos269/ResidueEscape.lean}{76}{no_bounded_positive_state_of_residue_escape}
\item \lref{Erdos269/ResidueEscape.lean}{96}{residue_le_bound_of_bounded_positive_state}
\item \lref{Erdos269/ResidueEscape.lean}{110}{no_bounded_positive_int_state_of_leastPositiveResidue}
\item \lref{Erdos269/RestrictedFloorSum.lean}{417}{windowBase}
\item \lref{Erdos269/RestrictedFloorSum.lean}{422}{windowForcing}
\item \lref{Erdos269/RestrictedFloorSum.lean}{455}{affineRecurrence_window}
\item \lref{Erdos269/RestrictedFloorSum.lean}{469}{windowForcing_const_mul}
\item \lref{Erdos269/RestrictedFloorSum.lean}{480}{integralCarry_window}
\item \lref{Erdos269/RestrictedFloorSum.lean}{548}{smooth3Val_dvd_threePrimeHeight_of_le}
\item \lref{Erdos269/RestrictedFloorSum.lean}{581}{integralCarry_cancel_commonFactor}
\item \lref{Erdos269/RestrictedFloorSum.lean}{601}{reducedCarry_pos_le_of_commonFactor}
\item \lref{Erdos269/RestrictedFloorSum.lean}{612}{reducedIntegralCarry_window}
\item \lref{Erdos269/RestrictedFloorSum.lean}{629}{CofinalLocalWindowEscape}
\item \lref{Erdos269/RestrictedFloorSum.lean}{645}{no_positive_reducedCarry_of_cofinalLocalWindowEscape}
\item \lref{Erdos269/ThreePrimeRunningLcm.lean}{721}{kernel_235_minor_eq_neg_one_fifteen}
\item \lref{Erdos269/ResidueEscape.lean}{138}{leastPositiveResidue_eq_natAbs_of_pos_le_modEq}
\item \lref{Erdos269/BoundedRadixTailEscape.lean}{89}{boundedRadix_zero_or_cofinal_far}
\item \lref{Erdos269/RestrictedFloorSum.lean}{497}{leastPositiveResidue_windowForcing_eq_carry}
\item \lref{Erdos269/WeightedPhaseCarry.lean}{109}{carry_eq_residueDigit_add_coboundary}
\item \lref{Erdos269/WeightedPhaseCarry.lean}{293}{CartierRealisation}
\item \lref{Erdos269/WeightedPhaseCarry.lean}{334}{finite_realisedSpan_of_factorisation}
\item \lref{Erdos269/WeightedPhaseCarry.lean}{150}{carryResidue_mem_interval}
\item \lref{Erdos269/WeightedPhaseCarry.lean}{157}{residueDigit_mem_interval}
\item \lref{Erdos269/ThreeChannelBlockRigidity.lean}{59}{channelBlockNull_iff_channelPotential}
\item \lref{Erdos269/CarryLiftExtinction.lean}{178}{no_carryLift_of_errorBound_below_twoPow}
\item \lref{Erdos269/CarryLiftExtinction.lean}{238}{no_bounded_carryLift_of_blockNull_twoAnchors}
\item \lref{Erdos269/CarryLiftExtinction.lean}{289}{binaryLift_twoAnchors_force_firstBlock_sum_one}
\item \lref{Erdos269/CarryLiftExtinction.lean}{308}{no_unitAccurateLift_with_twoAnchors_and_firstTwoBlockNull}
\end{itemize}
% END pin-faithful-source-inventory

\begin{thebibliography}{99}
\bibitem{erdosgraham1980} Paul Erd\H{o}s and Ronald L. Graham,
  \emph{Old and New Problems and Results in Combinatorial Number Theory}, Monographies de L'Enseignement Math\'{e}matique \textbf{28}, L'Enseignement Math\'{e}matique (1980),
  \href{https://mathweb.ucsd.edu/~ronspubs/80_11_number_theory.pdf}{source}.
\bibitem{erdos1988} Paul Erd\H{o}s,
  \emph{On the irrationality of certain series: problems and results}, in \emph{New Advances in Transcendence Theory}, Cambridge University Press (1988), 102--109,
  \href{https://doi.org/10.1017/CBO9780511897184.009}{doi:\nolinkurl{10.1017/CBO9780511897184.009}}.
\bibitem{erdos1974letter} Paul Erd\H{o}s,
  \emph{Letter to the Editor}, Fibonacci Quarterly \textbf{12}, no.~4 (1974), 335,
  \href{https://www.fq.math.ca/Scanned/12-4/letter.pdf}{source}.
\bibitem{erdosproblems} Thomas F. Bloom,
  \emph{Erd\H{o}s Problem \#269} (2026),
  \href{https://www.erdosproblems.com/269}{source}. Catalogue snapshot cited in the supplied manuscript: 28 July 2026.
\bibitem{bugeaudlaurent2023} Yann Bugeaud and Michel Laurent,
  \emph{Transcendence and continued fraction expansion of values of Hecke--Mahler series}, Acta Arithmetica \textbf{209} (2023), 59--90,
  \href{https://doi.org/10.4064/aa220323-18-1}{doi:\nolinkurl{10.4064/aa220323-18-1}}; arXiv:\href{https://arxiv.org/abs/2203.12901}{2203.12901}.
\bibitem{loxtonvdp1977} John H. Loxton and Alfred J. van der Poorten,
  \emph{Arithmetic properties of certain functions in several variables III}, Bulletin of the Australian Mathematical Society \textbf{16} (1977), 15--47,
  \href{https://doi.org/10.1017/S0004972700022978}{doi:\nolinkurl{10.1017/S0004972700022978}}.
\bibitem{fan2026comment} Steve Fan,
  \emph{Comment on Erd\H{o}s Problem \#269, thread 269, post 7218} (2026),
  \href{https://www.erdosproblems.com/forum/thread/269#post-7218}{source}. 26 June 2026, thread 269, post 7218; priority retained from the supplied record.
\bibitem{erdosstraus1974} Paul Erd\H{o}s and Ernst G. Straus,
  \emph{On the irrationality of certain series}, Pacific Journal of Mathematics \textbf{55}, no.~1 (1974), 85--92,
  \href{https://doi.org/10.2140/pjm.1974.55.85}{doi:\nolinkurl{10.2140/pjm.1974.55.85}}.
\bibitem{hancltijdeman2004} Jaroslav Han\v{c}l and Robert Tijdeman,
  \emph{On the irrationality of Cantor and Ahmes series}, Publicationes Mathematicae Debrecen \textbf{65}, no.~3--4 (2004), 371--380,
  \href{https://doi.org/10.5486/PMD.2004.3254}{doi:\nolinkurl{10.5486/PMD.2004.3254}}.
\bibitem{erdostaylor1957} Paul Erd\H{o}s and S. James Taylor,
  \emph{On the set of points of convergence of a lacunary trigonometric series and the equidistribution properties of related sequences}, Proceedings of the London Mathematical Society \textbf{s3-7}, no.~1 (1957), 598--615,
  \href{https://doi.org/10.1112/plms/s3-7.1.598}{doi:\nolinkurl{10.1112/plms/s3-7.1.598}}.
\bibitem{fan2026strongly} Steve Fan,
  \emph{Strongly complete sets and a conjecture of Erd\H{o}s} (2026),
  \href{https://arxiv.org/abs/2607.14071v1}{source}; arXiv:\href{https://arxiv.org/abs/2607.14071}{2607.14071}. The cited Lemma~3.1 is in arXiv v1, 15 July 2026.
\bibitem{hancltijdeman2008} Jaroslav Han\v{c}l and Robert Tijdeman,
  \emph{On the irrationality of polynomial Cantor series}, Acta Arithmetica \textbf{133}, no.~1 (2008), 37--52,
  \href{https://doi.org/10.4064/aa133-1-3}{doi:\nolinkurl{10.4064/aa133-1-3}}.
\bibitem{lucaouaknineworrell2025} Florian Luca, Jo\"{e}l Ouaknine and James Worrell,
  \emph{Transcendence of Hecke--Mahler Series}, Bulletin of the London Mathematical Society \textbf{57}, no.~5 (2025), 1360--1368,
  \href{https://doi.org/10.1112/blms.70033}{doi:\nolinkurl{10.1112/blms.70033}}; arXiv:\href{https://arxiv.org/abs/2412.07908}{2412.07908}. Numbered references use the published article.
\bibitem{kebis2024echoing} Pavol Kebis, Florian Luca, Jo\"{e}l Ouaknine, Andrew Scoones and James Worrell,
  \emph{On Transcendence of Numbers Related to Sturmian and Arnoux-Rauzy Words}, in \emph{51st International Colloquium on Automata, Languages, and Programming (ICALP 2024)}, Leibniz International Proceedings in Informatics \textbf{297}, Schloss Dagstuhl -- Leibniz-Zentrum f\"{u}r Informatik (2024), 144:1--144:15,
  \href{https://doi.org/10.4230/LIPIcs.ICALP.2024.144}{doi:\nolinkurl{10.4230/LIPIcs.ICALP.2024.144}}.
\bibitem{tijdemanmeijer1974} Robert Tijdeman and H. G. Meijer,
  \emph{On integers generated by a finite number of fixed primes}, Compositio Mathematica \textbf{29}, no.~3 (1974), 273--286,
  \href{https://www.numdam.org/article/CM_1974__29_3_273_0.pdf}{source}.
\bibitem{languasco2025} Alessandro Languasco, Florian Luca, Pieter Moree and Alain Togb\'{e},
  \emph{Sequences of integers generated by two fixed primes}, Abhandlungen aus dem Mathematischen Seminar der Universität Hamburg \textbf{95} (2025), 123--148,
  \href{https://doi.org/10.1007/s12188-025-00293-9}{doi:\nolinkurl{10.1007/s12188-025-00293-9}}; arXiv:\href{https://arxiv.org/abs/2309.12806}{2309.12806}.
\bibitem{kovactao2024} Vjekoslav Kova\v{c} and Terence Tao,
  \emph{On several irrationality problems for Ahmes series}, Acta Mathematica Hungarica \textbf{175} (2025), 572--608,
  \href{https://doi.org/10.1007/s10474-025-01528-0}{doi:\nolinkurl{10.1007/s10474-025-01528-0}}; arXiv:\href{https://arxiv.org/abs/2406.17593}{2406.17593}.
\bibitem{afperdosstraus2020} Angeliki Koutsoukou-Argyraki and Wenda Li,
  \emph{Irrationality Criteria for Series by Erd\H{o}s and Straus}, Archive of Formal Proofs (2020),
  \href{https://isa-afp.org/entries/Irrational_Series_Erdos_Straus.html}{source}. Entry dated 12 May 2020; proof-document version consulted: 6 February 2026.
\bibitem{adamczewskibugeaud2007} Boris Adamczewski and Yann Bugeaud,
  \emph{On the complexity of algebraic numbers I. Expansions in integer bases},
  Annals of Mathematics \textbf{165}, no.~2 (2007), 547--565,
  \href{https://doi.org/10.4007/annals.2007.165.547}{doi:\nolinkurl{10.4007/annals.2007.165.547}}.
\bibitem{adamczewskifaverjon2026} Boris Adamczewski and Colin Faverjon,
  \emph{Mahler's method in several variables and finite automata},
  Annals of Mathematics \textbf{204}, no.~2 (2026), 455--533,
  \href{https://doi.org/10.4007/annals.2026.204.2.1}{doi:\nolinkurl{10.4007/annals.2026.204.2.1}}.
  Online 13 September 2026; locators here refer to the
  \href{https://faverjon.perso.math.cnrs.fr/AdamczewskiFaverjon_MahlerFiniteAutomata.pdf}{68-page author manuscript}.
\end{thebibliography}

\companionpapercontext

\end{document}
